\section{Counting rational curves and obtaining the shell}
\label{sec:conclusion}

We apply the preceding analysis to every datum of
Lemma~\ref{compact:classes}, using the decay rate fixed after
Lemma~\ref{curve:area}. Proposition~\ref{compact:section} gives a divisor
$D_i$ with ordinary class $j(e_i)$. By
Theorem~\ref{cylinder:compactness}, every irreducible component of
$D_i$ is compact and has rational normalization. The divisor as a whole
may still have infinitely many components; it is locally finite and
bounded above in height. This distinction will be retained in all
intersection pairings below.

\begin{lemma}[Lifts of rational curves]\label{conclusion:lifts}
Every compact rational curve in $Y$ projects to a rational curve in $X$.
For each rational curve $B\subset X$, the compact irreducible curves in
$Y$ projecting onto $B$ form one orbit under $T$. In particular, a
collection of $r$ rational curves in $X$ gives at most $nr$ orbits of
compact lifts under $S=T^n$.
\end{lemma}

\begin{proof}
The image under $\pi$ of a compact curve is a compact analytic curve,
because the map restricted to that curve is proper and locally finite.
It is nonconstant since $\pi$ is a local biholomorphism. If the upstairs
normalization is $\PP^1$, the induced nonconstant map to the normalization
of the image is a finite map from $\PP^1$ to a compact Riemann surface.
Riemann--Hurwitz forces that surface to have genus zero. Thus the image
is rational.

Conversely, let $\nu:\PP^1\to B\subset X$ be the normalization of a
rational curve. Since $\PP^1$ is simply connected, the map $\nu$ lifts
to a holomorphic map $\widetilde\nu:\PP^1\to Y$ after a choice of the
image of one point. Its image $\widetilde B$ is compact and irreducible.
Over the smooth part of $B$, its projection has degree one, because
$\nu$ does. The lifts of this fixed normalization map form one orbit
under the deck group, so their images do also.

To see that these are all the compact irreducible curves over $B$, take
such a curve $C$ and a point of it over the smooth part of $B$.
Choose the lift of $\nu$ through this point. The local biholomorphism
property of $\pi$ makes its image and $C$ agree in a nonempty open
curve germ. Irreducibility and analytic continuation then make their
images equal. This proves the orbit assertion. Splitting the integer
deck powers into their $n$ residue classes proves the bound for $S$.

This argument lifts the normalization map, not the singular topological
space $B$. In particular it also applies to a nodal rational curve;
the connected components of its full inverse image need not be compact.
\end{proof}

One curve downstairs accounts for infinitely many translates upstairs.
We therefore count independent compact classes over a ring in which
multiplication by $t$ records one deck translation. A single orbit of
compact curves will then impose one linear equation over the fraction
field of that ring. The intersection pairing will show that fewer curve
orbits than independent classes are impossible.

Fix $n$, put $S=T^n$, and set
\[
 \Lambda=\mathbb Q[t,t^{-1}],\qquad
 \mathbb F=\mathbb Q(t),\qquad \overline{f(t)}=f(t^{-1}).
\]
On compactly supported cohomology, let
\[
 ta=S_!a:=(S^{-1})^*a.
\]
This is the convention that translates the support in the positive
$S$-direction. It agrees with the pushforward action $C\mapsto S_*C$
on homology under Poincar\'e duality. Define
\begin{equation}\label{conclusion:laurent-pairing}
 P(a,b)=\sum_{k\in\ZZ}\langle a,S_!^k b\rangle t^k
 \quad(a,b\in H_c^2(Y;\mathbb Q)).
\end{equation}
This is a Laurent polynomial: compactly supported representatives
intersect only finitely many translates of one another. Invariance of
intersection under $S$ and symmetry in degree two give
\begin{equation}\label{conclusion:sesquilinearity}
 P(ta,b)=tP(a,b),\qquad P(a,tb)=t^{-1}P(a,b),\qquad
 P(b,a)=\overline{P(a,b)}.
\end{equation}
Thus $P$ is linear in its first variable and conjugate-linear in its
second, where conjugation means only the involution $t\mapsto t^{-1}$.
For a compact oriented curve $C$, define similarly
\begin{equation}\label{conclusion:curve-pairing}
 Q(a,C)=\sum_{k\in\ZZ}\left(\int_{S_*^k C}a\right)t^k.
\end{equation}
This is finite and $\Lambda$-linear in $a$ with the same convention.
Its vanishing is equivalent to the vanishing of the ordinary pairing
of $a$ with every $S$-translate of $C$.

\begin{lemma}[Nonsingularity of the compact classes]
\label{conclusion:pairing-rank}
For the classes $e_1,\ldots,e_d$ of Lemma~\ref{compact:classes}, the
matrix $P(t)=(P(e_i,e_j))_{i,j=1}^d$ has nonzero determinant in
$\Lambda$. Their $\Lambda$-span is free of rank $d$, and $P$ is
nonsingular on its extension of scalars to $\mathbb F$.
\end{lemma}

\begin{proof}
Let $w_i$ be the transferred classes in $X_n$. Periodizing compactly
supported closed representatives of $e_i$ and $e_j$, then integrating
over one fundamental period, gives
\begin{equation}\label{conclusion:evaluation-one}
 P(e_i,e_j)(1)
 =\sum_{k\in\ZZ}\langle e_i,S_!^k e_j\rangle
 =w_i\cdot w_j.
\end{equation}
For completeness, the periodization is the locally finite sum of all
$S_!$-translates of a representative; it descends to a closed form on
$X_n$. In the integral of the wedge product of two such sums, change
variables by the translate in the first factor. The translated
fundamental periods tile $Y$, leaving the displayed sum indexed by the
relative translate of the second factor. Compact supports justify
every interchange. In the construction of Lemma~\ref{compact:classes}
the forms are supported inside the same open period, so terms with
$k\neq0$ actually vanish; the periodization identity also records the
intrinsic pairing convention.

The classes $w_i$ are independent over $\mathbb Q$ and the form on
$X_n$ is negative definite. Their Gram matrix is therefore negative
definite, in particular invertible. Equation~\eqref{conclusion:evaluation-one}
implies $\det P(1)\neq0$, so $\det P(t)$ is not the zero Laurent
polynomial. The matrix is invertible over $\mathbb F$.
If $\sum_i f_i(t)e_i=0$ for $f_i\in\Lambda$, pairing with each $e_j$
and using \eqref{conclusion:sesquilinearity} gives
$(f_1,\ldots,f_d)P(t)=0$. Invertibility over $\mathbb F$ forces every
$f_i=0$. This proves freeness and the asserted nonsingularity.
\end{proof}

\begin{proposition}[The required number of rational curves]
\label{conclusion:curves}
The surface $X$ contains at least $b_2(X)$ rational curves.
\end{proposition}

\begin{proof}
If there are infinitely many, the conclusion is immediate. Otherwise
suppose for a contradiction that their number is $r<b_2(X)$, and list
all of them. Choose $n$ so large that
\[
 n\bigl(b_2(X)-r\bigr)>C_0.
\]
Lemma~\ref{compact:classes} then gives $d>nr$ compact classes.
By Lemma~\ref{conclusion:lifts}, choose representatives
$C_1,\ldots,C_q$ for all $S$-orbits of compact irreducible lifts of the
listed rational curves, where $q\leq nr<d$.

Consider over $\mathbb F$ the $q$ homogeneous linear equations
\begin{equation}\label{conclusion:annihilating-equations}
 \sum_{i=1}^d f_i(t)\,Q(e_i,C_\ell)=0,
 \qquad 1\leq\ell\leq q.
\end{equation}
There are more unknowns than equations, so a nonzero solution exists.
Multiplying all its entries by a common nonzero denominator gives
$f_i\in\Lambda$, not all zero. The compactly supported class
\[
 x=\sum_{i=1}^d f_i(t)e_i
\]
is nonzero by Lemma~\ref{conclusion:pairing-rank}; its support is compact
because this is a finite Laurent combination. Equations
\eqref{conclusion:annihilating-equations} and
\eqref{conclusion:curve-pairing} imply that $x$ pairs to zero with
every compact rational lift and every one of its $S$-translates.

For each $j$, the ordinary class $j(e_j)$ is represented by the divisor
$D_j$. Every component of $D_j$ is compact rational by
Theorem~\ref{cylinder:compactness}, and its image downstairs is rational
by Lemma~\ref{conclusion:lifts}. Hence each component occurs among the
curves just annihilated by $x$. If $\alpha$ is a compactly supported
closed real two-form representing $x$, the ordinary divisor identity
\eqref{compact:divisor-pairing} gives, for every $k\in\ZZ$,
\[
 \langle x,S_!^k e_j\rangle
 =\int_{S_*^kD_j}\alpha=0.
\]
Indeed local finiteness makes the last integral a finite sum of
integrals over compact components meeting $\supp\alpha$, and each
summand is zero by the curve orthogonality. This argument pairs a
compact test class with a locally finite divisor; it never pairs two
locally finite cycles or sums an infinite Laurent series.

We have proved $P(x,e_j)=0$ for all $j$. By
\eqref{conclusion:sesquilinearity} this says
$(f_1,\ldots,f_d)P(t)=0$, contradicting the nonsingularity of $P(t)$
over $\mathbb F$. Therefore $r\geq b_2(X)$.
\end{proof}

\begin{lemma}[Independence of curve classes]
\label{conclusion:upper-bound}
Under \eqref{surface:no-square-zero}, the real cohomology classes of
any finite collection of distinct curves on $X$ are linearly
independent. Consequently $X$ has at most $b_2(X)$ curves in total.
\end{lemma}

\begin{proof}
Suppose a finite collection admits a nonzero real relation. Split its
positive and negative coefficients to obtain effective real divisors
$A$ and $B$ with disjoint lists of components and $[A]=[B]=v$; one of
them may be empty. If $v\neq0$, both are nonempty and negative
definiteness gives $v^2<0$. But positivity of local intersections of
distinct complex curves gives
\[
 v^2=A\cdot B\geq0,
\]
a contradiction. Thus $v=0$. At least one nonempty effective real
divisor, say $A=\sum_{i\in I}a_i B_i$ with all $a_i>0$, has zero real
class.

To pass from this real relation to an integral one, express the classes
$[B_i]$ in an integral basis of $H^2(X;\ZZ)/\operatorname{Tors}$.
The relation space is the kernel of an integer matrix, hence has a
basis over $\mathbb Q$. Its rational points are dense in its real
points. Approximate $(a_i)_{i\in I}$ within this kernel by a rational
vector with all coordinates still positive, and clear denominators.
The result is a nonzero effective integral divisor with zero real
class and hence square zero. This contradicts
\eqref{surface:no-square-zero}. There is no relation. Since
$\dim H^2(X;\RR)=b_2(X)$, any collection of more than $b_2(X)$ curves
would contain such a finite dependent collection.
\end{proof}

\begin{proof}[Proof of Theorem~\ref{thm:main}]
If $X$ has a nonzero effective square-zero divisor,
Proposition~\ref{surface:reduction}\textup{(i)} already gives the
desired shell. In the remaining case,
Proposition~\ref{conclusion:curves} and
Lemma~\ref{conclusion:upper-bound} show that $X$ has exactly $b_2(X)$
rational curves. It is compact and minimal, has $b_1(X)=1$,
$\kappa(X)=-\infty$, and $b_2(X)>0$ by the original hypotheses.
Every hypothesis of the Main Theorem of Dloussky--Oeljeklaus--Toma is
therefore satisfied, so Proposition~\ref{surface:reduction}\textup{(ii)}
applies.

The conclusion in \cite[pp.~283--284]{DOT2003} is an open neighborhood
$U$ of the unit sphere in $\CC^2\setminus\{0\}$ and a biholomorphism
$\phi:U\to\Sigma\subset X$ such that $X\setminus\Sigma$ is connected.
This is precisely the conclusion required in Theorem~\ref{thm:main}.
Equivalently, if a shell is described by its nonseparating core sphere,
choose a sufficiently small round annular collar in the given
holomorphic neighborhood. On either side of the core, the collar
coordinate pushes the complement of the core onto the complement of
the smaller open collar. The latter is therefore connected as well.
\end{proof}

\subsection{Consequences of the shell}\label{sec:consequences}

The shell now gives the deformation and topological description stated
in the introduction, and removes the remaining class VII exception in
the aspherical-surface inequality.

\begin{proof}[Proof of Corollary~\ref{cor:topology}]
Theorem~\ref{thm:main} supplies a global spherical shell. Kato's
deformation theorem \cite[Theorem~1]{Kato1977} then gives a
one-parameter deformation with each nonzero fiber a modification of a
primary Hopf surface; his Corollary~1 gives $\pi_1(X)\cong\ZZ$.
In complex dimension two, a modification between smooth surfaces factors
as a finite succession of point blowups, with possibly infinitely near
centers \cite{BHPV2004}.

After shrinking $\Delta$, the family is a proper smooth submersion.
Ehresmann's theorem identifies its fibers by diffeomorphisms preserving
their continuously varying complex orientations. Every primary Hopf
surface is diffeomorphic to $S^1\times S^3$
\cite[\S1, p.~1109]{GauduchonOrnea1998}, and hence has second Betti number
zero. Each complex point blowup increases the second Betti number by one
and replaces the oriented smooth manifold by its connected sum with
$\overline{\mathbb{CP}}^{\,2}$ \cite{BHPV2004}. Since the fibers have
second Betti number $b$, there are exactly $b$ blowups, and the asserted
oriented smooth type follows. The product orientation can be chosen
compatibly because $S^1\times S^3$ admits an orientation-reversing
diffeomorphism.
\end{proof}

\begin{proof}[Proof of Corollary~\ref{cor:aspherical-geography}]
Albanese--Di Cerbo--Lombardi~\cite[Theorem~1.4]{AlbaneseDiCerboLombardi2025}
prove these conclusions except for the possible case of a class
$\mathrm{VII}_0^+$ surface violating the Global Spherical Shell conjecture.
Theorem~\ref{thm:main} excludes that case. Corollary~\ref{cor:topology}
also explains the non-asphericity: such a surface has $\pi_1\cong\ZZ$
and $b_2>0$, but an aspherical manifold with fundamental group $\ZZ$
is homotopy equivalent to $S^1$.
\end{proof}
